\chapter{Energy, life and functional balance}\label{ch:homeostasis}

\section{The rule that prevents a free lunch}
A pump lifts water because something supplies energy. Electricity reaches its motor through a network; that electricity was generated from an identified source. The pump cannot deliver indefinitely from an empty energy supply, however persuasive a contract or a price might be. This is the ordinary force of the first law of thermodynamics.

For a closed thermodynamic system under a familiar sign convention,
\begin{equation}
 \Delta U=Q-W ,
\end{equation}
where $U$ is internal energy, $Q$ is heat entering and $W$ is work done by the system. An open system also carries energy with matter crossing its boundary. The equation does not say that every useful form of energy remains available for reuse; it says that a claimed change requires a corresponding transfer or transformation. It is the scientific reason a perpetual-motion machine of the first kind cannot be built.\cite{intro-energy}

The bookkeeping lesson is immediate. If a pump moves water while the account contains no energy source, either the physical description is incomplete or the pump has been granted impossible power. Recording the electricity as a flow does not decide whether pumping is wise. It prevents us from calling the action self-sustaining when it depends on something outside the picture.

\section{Why conserved energy is not enough}
The second law addresses a different question: which transformations are possible and what happens to the ability of energy to do useful work. For an isolated system, total entropy does not decrease. A real heat engine cannot convert every unit of heat from one reservoir into work while leaving no other change. Friction and mixing can preserve energy while dispersing it in forms less useful for the task at hand.\cite{intro-energy}

The Earth is not an isolated box. It receives energy from the Sun and radiates energy to space. A living organism can build and maintain local order by taking in energy and matter and exporting heat and waste. The second law does not forbid life; it explains why maintenance has a physical cost and why an account that shows only the organised result is incomplete.

At the clinic, a refrigerator preserves medicine by moving heat from its interior to the surrounding room while consuming electricity. Saying that the medicine remained cold without recording the power and the rejected heat would mistake a continuing service for a free property of the refrigerator. A durable economy likewise has to ask what flows sustain its useful arrangements.

\section{The body does not maximise its temperature}
Now consider a human body. Its temperature, blood pressure and concentrations of important substances vary while the body remains alive. Regulation does not mean setting every variable to zero, holding every amount perfectly fixed or making one quantity as large as possible. Too little and too much can both be dangerous. Physiology describes homeostasis as the maintenance of conditions compatible with function through interacting processes of sensing and response. Negative feedback is a central pattern: a departure triggers a response that tends to oppose it.\cite{intro-homeostasis,intro-physiology}

The point is clearer in a simple tank. Water enters and leaves, yet its level can remain within a useful region. If outflow increases, a sensor may detect a fall and a valve may admit more water. If the inlet fails, the level continues to fall despite the sensor. If the response is delayed, it may overshoot. The difference between a measurement, an available action and a successful response matters as much in a body as in an infrastructure network.

We sometimes speak as if growth or accumulation were an adequate objective for every system. The body shows why that cannot be a general physical rule. More heat is not always better; more glucose is not always better; more pressure is not always better. The same warning applies to a society that treats increased throughput or income as sufficient evidence of continuing health. A gain has to be assessed in relation to the conditions that make the system function.

\section{A balance that remains active}
Suppose the clinic's water tank contains ten units at the beginning and end of an hour. During the hour it receives four and supplies four. Another tank also begins and ends at ten but has both valves closed. Their equal snapshots conceal different service histories. This is why functional balance is not simply equality of two numbers: it can be a maintained process.

The stock account for the first tank is
\begin{equation}
 x_{t+1}=x_t+q_{\mathrm{in}}-q_{\mathrm{out}},
\end{equation}
with each $q$ measured in the same physical unit over the same interval. A constant stock means $q_{\mathrm{in}}=q_{\mathrm{out}}$ for that interval; it does not tell us whether the water was safe or whether a patient received it. Those are additional facts. At a larger scale, regeneration, consumption and disturbance have to be represented if we want a claim about continuing function.

\begin{intuitionbox}{Three separate tests}
Can the movement physically occur? Does the represented state change in a desirable direction? Can the process continue under demand and disturbance? The three questions belong together, but no one answer proves the other two.
\end{intuitionbox}

\section{Where should the system return?}
Negative feedback is intelligible only relative to a condition from which the system has departed. But what is that condition? A fixed number may be appropriate for one regulated variable and inappropriate for another. A body is not healthy because every variable is equal, and two towns need not hold equal stores. Nor does a mathematical minimum automatically cause physical motion toward itself.

We now need to separate a \emph{reference} used for comparison from an \emph{equilibrium} of a stated law of change and from the broader idea of \emph{homeostasis} through time. Those distinctions will let us say what it would mean for an action to restore rather than merely increase something. The next chapter builds them before any new accounting unit is introduced.
